\documentclass{MO}

\author{Dominik Dembinný}
\studentclass{2. F / 4 Gymnázium}
\school{Arabská 14, Praha 6}

\problemcategory{75}{B}{I}{4}

\begin{document}
\makeMOheader

\section*{Pozorování}
\subsection*{Definiční obor}
Budeme pozorovat, jakých hodnot může nabývat $a$ a $b$, tak aby splňovali následující nerovnice:
\begin{align}
    a + b &> 0 \\
    ab &> 0
\end{align}
\begin{center}
    \begin{tabular}{|c|c|c|c|c|}
        \hline
        $a$ & $b$ & $ab$ & $a + b$ & (1) $\land$ (2) \\ \hline
        $+$ & $+$ & $+$ & $+$ & ano \\ \hline
        $-$ & $+$ & $-$ & $?$ & ne \\ \hline
        $-$ & $-$ & $+$ & $-$ & ne \\ \hline
        $0$ & $+$ & $0$ & $+$ & ne \\ \hline
    \end{tabular}
\end{center}
Z tabulky můžeme vyčíst, že $a, b \in \mathbb{R}^{+}$.

\section*{Řešení}
Nejmenší hodnota a + b bude ležet na hranici, kde $ab = a + b$.
Nejprve najdeme pomocí derivace sklon $\dv{b}{a}$ v minimu.
\begin{align*}
    S &= a + b \\
    dS &= da + db \\
    da + db &= 0 \\
    \dv{b}{a} &= -1
\end{align*}
\begin{align*}
    ab &= a + b \\
    \dv{a} (ab) &= \dv{a} (a + b) \\
    b + a \dv{b}{a} &= 1 + \dv{b}{a} \\
    b + a \cdot (-1) &= 1 + (-1) \\
    b - a &= 0 \\
    &\implies a = b
\end{align*}
Dokázali jsme, že $a$ se musí rovnat v minimu $b$.
\begin{align*}
    a^2 &= 2a \\
    a > 0 &\implies a = 2 \land b = 2 \\
    S &= 4
\end{align*}
Součet $a + b$ musí být nejméně $4$, aby nerovnosti platily.

\end{document}